\section*{第 \ref{ch:b3:climate-models} 章——气候模式}

\begin{solution}{exo:b3:climate-models:1}
$A = 14.0 - 1.5 = \num{12.5}$ 且 $B = 0.14 - 0.05 = \num{0.09}$，单位为
\unit{kcal.cm^{-2}} 每月。一个月为 \qty{2.63e6}{s}，故
\qty{1}{kcal.cm^{-2}} 每月等于 $\num{4.184e7}/\num{2.63e6} =
\qty{15.91}{W/m^2}$：$A = \qty{199}{W/m^2}$，$B =
\qty{1.43}{W.m^{-2}.\celsius^{-1}}$。在 \qty{15}{\celsius} 时，$199 + 21.5 =
\qty{220}{W/m^2}$。
\end{solution}

\begin{solution}{exo:b3:climate-models:2}
$(238.0 - 211.9)/1.816 = \qty{14.4}{\celsius}$；取 0.31 时，$(234.6 -
211.9)/1.816 = \qty{12.5}{\celsius}$。反照率每增加 0.01，相当于
\qty{3.4}{W/m^2}，而 $3.4/1.816 = \qty{1.9}{\celsius}$。
\end{solution}

\begin{solution}{exo:b3:climate-models:3}
$C = 1025\times3992\times60 = \qty{2.46e8}{J.m^{-2}.K^{-1}}$；除以
一年的 \qty{3.156e7}{s}，得 \qty{7.78}{W.\text{年}.m^{-2}.K^{-1}}。于是 $\tau
= 7.78/1.816 = \num{4.3}\text{ 年}$，$7.78/1.33 = \num{5.8}\text{ 年}$。
\end{solution}

\begin{solution}{exo:b3:climate-models:4}
ECS $= 3.93/1.33 = \qty{2.95}{\celsius}$；TCR $\approx 3.93/1.95 =
\qty{2.02}{\celsius}$；比值为 \num{0.68}。
\end{solution}

\begin{solution}{exo:b3:climate-models:5}
$Q = \qty{335}{W/m^2}$。雪球态：$(335\times0.38 - 211.9)/1.816 =
\qty{-46.6}{\celsius}$；暖态：$(335\times0.70 - 211.9)/1.816 =
\qty{12.4}{\celsius}$，仍高于 \qty{10}{\celsius}；斜坡：$335(0.54 +
0.016\,T) = 211.9 + 1.816\,T$，解得 $T = \qty{8.7}{\celsius}$。$N' =
\qty{-1.816}{W.m^{-2}.\celsius^{-1}}$，前两个平衡点稳定；而 $-1.816
+ 335\times0.016 = \num{+3.54}$，斜坡处不稳定。
\end{solution}

\begin{solution}{exo:b3:climate-models:6}
暖态分支止于其温度达到 \qty{0}{\celsius} 之处：
$0.70\,Q = A$，$Q = \qty{302.7}{W/m^2}$，对应太阳常数为
\qty{1211}{W/m^2}，比今天低 \qty{11.0}{\percent}。雪球融化的条件为
$0.38\,Q = A - 20B$，$Q = \qty{462.1}{W/m^2}$，$S_0 =
\qty{1848}{W/m^2}$。阈值为 $B(T^* - T_F)/(Q_0(1 - \alpha_F))$，
即今天的温度 $T^*$ 与冰开始扩展时的温度之差，
再换算为太阳辐射：$1.816\times14.37/238 =
\qty{11.0}{\percent}$（此处），正文中则为
$1.816\times4.37/238 = \qty{3.3}{\percent}$。
\end{solution}

\begin{solution}{exo:b3:climate-models:7}
(a)~\qty{-45.5}{\celsius}。(b)~$340\times0.38 + F = 211.9 - 18.2$：$F =
\qty{64.5}{W/m^2}$，为一次加倍强迫的 16 倍。(c)~强迫持续期间，
行星跃迁至暖态分支，$(238.0 + 64.5 -
211.9)/1.816 = \qty{50}{\celsius}$；强迫撤除后，行星沿暖态分支冷却
至 \qty{14.4}{\celsius}，该温度在今天的太阳辐照下存在：
它不会回到雪球态。
\end{solution}

\begin{solution}{exo:b3:climate-models:8}
$\Sigma = 1.80/7.0 + 0.70/90 = \qty{0.2649}{\text{年}^{-1}}$，$\Pi = 1.10\times0.70/
(7.0\times90) = \qty{1.222e-3}{\text{年}^{-2}}$。由根可得 $\tau_f =
\num{3.84}\text{ 年}$，$\tau_s = \num{213}\text{ 年}$；其近似值为
$7.0/1.80 = \num{3.89}\text{ 年}$，$90\times1.80/(1.10\times0.70) =
\num{210}\text{ 年}$。$a_f = 0.597$ 精确，$1.10/1.80 = 0.611$
为近似值。ECS $= 3.93/1.10 = \qty{3.57}{\celsius}$；TCR $\approx
3.93/1.80 = \qty{2.18}{\celsius}$。
\end{solution}

\begin{solution}{exo:b3:climate-models:9}
分母 $B + \beta = \num{5.555}$。赤道：$(340\times1.221\times0.70 -
211.9 + 3.74\times12.81)/5.555 = \qty{22.8}{\celsius}$；极区：
$(340\times0.506\times0.38 - 211.9 + 47.9)/5.555 = \qty{-17.7}{\celsius}$。
输送项在赤道处移除 $3.74\times(22.8 - 12.8) = \qty{37}{W/m^2}$，
并向极区输送 $3.74\times30.6 = \qty{114}{W/m^2}$。
\end{solution}

\begin{solution}{exo:b3:climate-models:10}
(a)~$(1 - x^2)P_2' = 3x - 3x^3$，其导数为 $3 - 9x^2 = -6P_2$。
(b)~输送项为 $-\mathcal D T_2(-6P_2) = 6\mathcal DT_2P_2$。
平衡式 $Q(1 + s_2P_2)(1 - \alpha) - A - B(T_0 + T_2P_2) - 6\mathcal
DT_2P_2 = 0$ 对所有 $x$ 成立，若其常数部分与 $P_2$ 部分
均为零，由此得 $T_0$ 与 $T_2$。(c)~$T(0) - T(1) = T_2(-\tfrac12 - 1)
= -1.5\,T_2$。若无输送，$T_2 = 340\times(-0.476)\times0.70/1.816 =
\qty{-62.4}{\celsius}$：差值为 \qty{94}{\celsius}。(d)~$T_2 =
-55.3/1.5 = \qty{-36.9}{\celsius}$，故 $B + 6\mathcal D = 113.3/36.9 =
3.07$，$\mathcal D = \qty{0.21}{W.m^{-2}.\celsius^{-1}}$。
\end{solution}

\begin{solution}{exo:b3:climate-models:11}
(a)~$r = \tfrac12\Sigma\big(1 \pm \sqrt{1 - 4\Pi/\Sigma^2}\big)$；取
$\sqrt{1 - \epsilon} \approx 1 - \epsilon/2$，两根为 $\Sigma -
\Pi/\Sigma \approx \Sigma$ 与 $\Pi/\Sigma$。(b)~令 $T_d = 0$，$C\,\mathrm
dT/\mathrm dt = kt - (\gamma - \lambda)T$，自静止出发的解为 $T =
k(t - \tau_f)/(\gamma - \lambda) + k\tau_f\mathrm e^{-t/\tau_f}/(\gamma -
\lambda)$，$\tau_f = C/(\gamma - \lambda)$。当 $t \gg \tau_f$ 时，$T \approx
kt/(\gamma - \lambda)$。在加倍时刻，$kt = F_{2\times}$：TCR
$\approx F_{2\times}/(\gamma - \lambda)$。(c)~$\Sigma = 1.95/8.1 +
0.62/110 = \qty{0.2464}{\text{年}^{-1}}$，$\Pi = 1.33\times0.62/(8.1\times110) =
\qty{9.25e-4}{\text{年}^{-2}}$；$\tau_f = \num{4.12}\text{ 年}$，$\tau_s =
\num{262}\text{ 年}$（近似为 $1/\Sigma = \num{4.06}\text{ 年}$，而
$\Sigma/\Pi = \num{266}\text{ 年}$）；$4\Pi/\Sigma^2 = 0.061$。
\end{solution}

\begin{solution}{exo:b3:climate-models:12}
(a)~均值 \qty{3.50}{\celsius}，标准差 \qty{0.81}{\celsius}。
(b)~ECS $= 1.88 + 1.86\,X$；相关系数 \num{0.992}；残差
标准差 \qty{0.11}{\celsius}。(c)~$1.88 + 1.86\times0.7 =
\qty{3.18}{\celsius}$，散布为 $\sqrt{(1.86\times0.1)^2 + 0.11^2} =
\qty{0.22}{\celsius}$，约为集合散布的四分之一。
(d)~$X$ 与反馈之间没有物理机制相联系；$X$ 是通过尝试大量可观测量而找到的，
因此相关性可能只是偶然；所有八个模式共有一种偏差；
同一系列的模式被当作
彼此独立。
\end{solution}

\begin{solution}{pb:b3:climate-models:1}
\textbf{1.} $Q_0 = 1360/4 = \qty{340}{W/m^2}$：行星以
面积为 $\pi R^2$ 的圆盘截获太阳光，再将其摊铺到面积为
$4\pi R^2$ 的球面上。

\textbf{2.} $340\times0.70 = \qty{238.0}{W/m^2}$；$211.9 + 1.816\times15 =
\qty{239.1}{W/m^2}$。

\textbf{3.} $(238.0 - 211.9)/1.816 = \qty{14.4}{\celsius}$。

\textbf{4.} $(129.2 - 211.9)/1.816 = \qty{-45.5}{\celsius}$，低于
\qty{-10}{\celsius}。

\textbf{5.} $340(0.54 + 0.016\,T) = 211.9 + 1.816\,T$：$T =
(183.6 - 211.9)/(1.816 - 5.44) = \qty{7.8}{\celsius}$。

\textbf{6.} $(238.0/\num{5.670e-8})^{1/4} = \qty{254.5}{K}$，
\qty{-18.6}{\celsius}：比暖平衡态低 \qty{33}{\celsius}，其中行星的温室效应
已包含在 $A$ 和 $B$ 之中。

\textbf{7.} 在雪球态与暖平衡态处，$N' = -B = \qty{-1.816}{W.m^{-2}.\celsius^{-1}}$，两者均稳定；
在中间的平衡态处，$-1.816 + 5.44 =
\num{+3.62}$，不稳定。

\textbf{8.} $\tau = 8.1/3.62 = \num{2.2}\text{ 年}$。

\textbf{9.} $8.1/1.816 = \num{4.5}\text{ 年}$。

\textbf{10.} $3.93/1.816 = \qty{2.2}{\celsius}$，低于评估的最佳估计值
\qty{3}{\celsius}，也低于其\emph{可能}范围，即 2.5 至
\qty{4}{\celsius}：模式的 $B$ 由跨纬度拟合得到，
并不等于整个气候的反馈参数。

\textbf{11.} $\mathrm e^{-t/\tau} = 0.05$，取 $t = 3\tau =
\num{13}\text{ 年}$。实际气候中深海同样在变暖，其时间
尺度为若干世纪。

\textbf{12.} 令 $T_F = \qty{10}{\celsius}$；$Q = (211.9 + 18.2)/0.70 =
\qty{328.7}{W/m^2}$。

\textbf{13.} $S_0 = \qty{1315}{W/m^2}$，比今天低 \qty{3.3}{\percent}。

\textbf{14.} 对雪球态，$(328.7\times0.38 - 211.9)/1.816 =
\qty{-47.9}{\celsius}$。

\textbf{15.} $Q = (211.9 - 18.2)/0.38 = \qty{509.8}{W/m^2}$，$S_0 =
\qty{2039}{W/m^2}$，比今天高 \qty{50}{\percent}。

\textbf{16.} 在 \qty{1315}{W/m^2} 下沿暖支降至 \qty{10}{\celsius}，
随后降至 \qty{-48}{\celsius}，然后沿雪球支回升至 \qty{-45.5}{\celsius}，
对应 \qty{1360}{W/m^2}：行星
仍保持冰冻。只有超过 \qty{2039}{W/m^2}，
它才会跳回暖支。

\textbf{17.} $\bar\alpha = 0.30\times0.9738 + 0.62\times0.0262 =
\num{0.3084}$；$\bar T = (340\times0.6916 - 211.9)/1.816 =
\qty{12.80}{\celsius}$。

\textbf{18.} $T_c = (340\times0.5643\times0.50 - 211.9 + 3.739\times12.80)/
5.555 = \qty{-12.26}{\celsius}$，比 \ang{72} 处的再分析结果
高 \qty{3.5}{\celsius}。

\textbf{19.} $5.555\times(-12.26) + 211.9\times3.059 = \num{580.1}$，其中
$\beta/B = \num{2.059}$。

\textbf{20.} 各分母依次为 1.7053、1.7473、1.7565、1.7488 和
1.7362，而 $Q/Q_0 = \num{1.0005}$、\num{0.9764}、\num{0.9713}、
\num{0.9756} 和 \num{0.9827}。

\textbf{21.} 极小值位于 \ang{47.5} 附近，为 \num{0.971}。极小值极向一侧的冰线
是稳定的，$\mathrm dQ/\mathrm dx_s > 0$；赤道向一侧的冰线
则不稳定。

\textbf{22.} 在 \ang{30} 与 \ang{40} 之间，$Q/Q_0 = 1$ 也成立，且接近
\ang{30}（插值得 \ang{30.2}）处为不稳定的冰线。在 \ang{0}，
$Q/Q_0 = 580.1/(0.6106 + 0.7824)/340 = \num{1.225}$：雪球态在比今天亮至多
\qty{22}{\percent} 的太阳下都能维持，
因此今天仍有可能出现雪球态。

\textbf{23.} \qty{2.9}{\percent}，对照零维模式的 \qty{3.3}{\percent}
与布迪科的 \qty{1.6}{\percent}；阈值处冰线位于
\ang{48} 附近，布迪科的结果为 \ang{50}。纬度模式
比全球模式更早冰封，因为日照在
冰形成的地方最弱。

\textbf{24.} 模式行星冰封所需的太阳常数降幅
为 \textbf{\qty[text-series-to-math=true]{2.9}{\percent}}
（按计算为 \qty{2.87}{\percent}）。
\end{solution}
